% Chunk: Introduction and Section 1.1, Relativistic Wave Mechanics.
% Source: Weinberg QFT Vol. I, physical PDF pages 24--37 (printed pages 1--14).
% Status: source-reviewed and compile-clean.
% Physical pages: 24--37 (printed pages 1--14).
% Figures: none.
% Uncertainties: none.

Our immersion in the present state of physics makes it hard for us to
understand the difficulties of physicists even a few years ago, or to profit
from their experience. At the same time, a knowledge of our history is a
mixed blessing---it can stand in the way of the logical reconstruction of
physical theory that seems to be continually necessary.

I have tried in this book to present the quantum theory of fields in
a logical manner, emphasizing the deductive trail that ascends from the
physical principles of special relativity and quantum mechanics. This
approach necessarily draws me away from the order in which the subject
in fact developed. To take one example, it is historically correct that
quantum field theory grew in part out of a study of relativistic wave
equations, including the Maxwell, Klein--Gordon, and Dirac equations.
For this reason it is natural that courses and treatises on quantum field
theory introduce these wave equations early, and give them great weight.
Nevertheless, it has long seemed to me that a much better starting point is
Wigner's definition of particles as representations of the Poincar\'e
group, even though this work was not published until 1939 and
did not have a great impact for many years after. In this book we start
with particles and get to the wave equations later.

This is not to say that particles are necessarily more fundamental than
fields. For many years after 1950 it was generally assumed that the
laws of nature take the form of a quantum theory of fields. I start
with particles in this book, not because they are more fundamental, but
because what we know about particles is more certain, more directly
derivable from the principles of quantum mechanics and relativity. If it
turned out that some physical system could not be described by a quantum
field theory, it would be a sensation; if it turned out that the system did
not obey the rules of quantum mechanics and relativity, it would be a
cataclysm.

In fact, lately there has been a reaction against looking at quantum
field theory as fundamental. The underlying theory might not be a theory
of fields or particles, but perhaps of something quite different, like strings.
From this point of view, quantum electrodynamics and the other quantum
field theories of which we are so proud are mere `effective field theories,'
low-energy approximations to a more fundamental theory. The reason that
our field theories work so well is not that they are fundamental truths,
but that any relativistic quantum theory will look like a field theory when
applied to particles at sufficiently low energy. On this basis, if we want to
know why quantum field theories are the way they are, we have to start
with particles.

But we do not want to pay the price of altogether forgetting our past.
This chapter will therefore present the history of quantum field theory
from earliest times to 1949, when it finally assumed its modern form. In
the remainder of the book I will try to keep history from intruding on
physics.

One problem that I found in writing this chapter is that the history of
quantum field theory is from the beginning inextricably entangled with
the history of quantum mechanics itself. Thus, the reader who is familiar
with the history of quantum mechanics may find some material here that
he or she already knows, especially in the first section, where I discuss the
early attempts to put together quantum mechanics with special relativity.
In this case I can only suggest that the reader should skip on to the less
familiar parts.

On the other hand, readers who have no prior familiarity with quantum
field theory may find parts of this chapter too brief to be altogether clear.
I urge such readers not to worry. This chapter is not intended as a
self-contained introduction to quantum field theory, and is not needed as
a basis for the rest of the book. Some readers may even prefer to start
with the next chapter, and come back to the history later. However, for
many readers the history of quantum field theory should serve as a good
introduction to quantum field theory itself.

I should add that this chapter is not intended as an original work
of historical scholarship. I have based it on books and articles by real
historians, plus some historical reminiscences and original physics articles
that I have read. Most of these are listed in the bibliography given at the
end of this chapter, and in the list of references. The reader who wants
to go more deeply into historical matters is urged to consult these listed
works.

A word about notation. In order to keep some of the flavor of past
times, in this chapter I will show explicit factors of $\hbar$ and $c$ (and even
$h$), but in order to facilitate comparison with modern physics literature,
I will use the more modern rationalized electrostatic units for charge, so
that the fine structure constant $\alpha\simeq 1/137$ is $e^2/4\pi\hbar c$. In subsequent
chapters I will mostly use the `natural' system of units, simply setting
$\hbar=c=1$.

\subsection{Relativistic Wave Mechanics}

Wave mechanics started out as relativistic wave mechanics. Indeed, as
we shall see, the founders of wave mechanics, Louis de Broglie and
Erwin Schr\"odinger, took a good deal of their inspiration from special
relativity. It was only later that it became generally clear that relativistic
wave mechanics, in the sense of a relativistic quantum theory of a fixed
number of particles, is an impossibility. Thus, despite its many successes,
relativistic wave mechanics was ultimately to give way to quantum field
theory. Nevertheless, relativistic wave mechanics survived as an important
element in the formal apparatus of quantum field theory, and it posed a
challenge to field theory, to reproduce its successes.

The possibility that material particles can like photons be described in
terms of waves was first suggested~\hyperref[ref:1.1]{[1]} in 1923 by Louis de Broglie. Apart
from the analogy with radiation, the chief clue was Lorentz invariance: if
particles are described by a wave whose phase at position $\mathbf{x}$ and time $t$
is of the form $2\pi(\boldsymbol{\kappa}\cdot\mathbf{x}-\nu t)$, and if this phase is to be Lorentz invariant,
then the vector $\boldsymbol{\kappa}$ and the frequency $\nu$ must transform like $\mathbf{x}$ and $t$, and
hence like $\mathbf{p}$ and $E$. In order for this to be possible $\boldsymbol{\kappa}$ and $\nu$ must have the
same velocity dependence as $\mathbf{p}$ and $E$, and therefore must be proportional
to them, with the same constant of proportionality. For photons, one had
the Einstein relation $E=h\nu$, so it was natural to assume that, for material
particles,
\begin{equation}
  \boldsymbol{\kappa}=\mathbf{p}/h,
  \qquad
  \nu=E/h,
  \tag{1.1.1}\label{eq:1.1.1}
\end{equation}
just as for photons. The group velocity $\partial\nu/\partial\boldsymbol{\kappa}$ of the wave then turns
out to equal the particle velocity, so wave packets just keep up with the
particle they represent.

By assuming that any closed orbit contains an integral number of
particle wavelengths $\lambda=1/|\boldsymbol{\kappa}|$, de Broglie was able to derive the old
quantization conditions of Niels Bohr and Arnold Sommerfeld, which
though quite mysterious had worked well in accounting for atomic spectra.
Also, both de Broglie and Walter Elsasser~\hyperref[ref:1.2]{[2]} suggested that de Broglie's
wave theory could be tested by looking for interference effects in the
scattering of electrons from crystals; such effects were established a few
years later by Clinton Joseph Davisson and Lester H. Germer~\hyperref[ref:1.3]{[3]}. However,
it was still unclear how the de Broglie relations \eqref{eq:1.1.1} should be modified
for non-free particles, as for instance for an electron in a general Coulomb
field.

Wave mechanics was by-passed in the next step in the history of
quantum mechanics, the development of matrix mechanics~\hyperref[ref:1.4]{[4]} by Werner
Heisenberg, Max Born, Pascual Jordan and Wolfgang Pauli in the years
1925--1926. At least part of the inspiration for matrix mechanics was the
insistence that the theory should involve only observables, such as the
energy levels, or emission and absorption rates. Heisenberg's 1925 paper
opens with the manifesto: `The present paper seeks to establish a basis
for theoretical quantum mechanics founded exclusively upon relationships
between quantities that in principle are observable.' This sort of positivism
was to reemerge at various times in the history of quantum field theory,
as for instance in the introduction of the S-matrix by John Wheeler and
Heisenberg (see Chapter 3) and in the revival of dispersion theory in the
1950s (see Chapter 10), though modern quantum field theory is very far
from this ideal. It would take us too far from our subject to describe
matrix mechanics in any detail here.

As everyone knows, wave mechanics was revived by Erwin Schr\"odinger.
In his 1926 series of papers~\hyperref[ref:1.5]{[5]}, the familiar non-relativistic wave equation
is suggested first, and then used to rederive the results of matrix mechanics.
Only later, in the sixth section of the fourth paper, is a relativistic
wave equation offered. According to Dirac~\hyperref[ref:1.6]{[6]}, the history is actually quite
different: Schr\"odinger first derived the relativistic equation, then became
discouraged because it gave the wrong fine structure for hydrogen, and
then some months later realized that the non-relativistic approximation
to his relativistic equation was of value even if the relativistic equation
itself was incorrect! By the time that Schr\"odinger came to publish his
relativistic wave equation, it had already been independently rediscovered
by Oskar Klein~\hyperref[ref:1.7]{[7]} and Walter Gordon~\hyperref[ref:1.8]{[8]}, and for this reason it is usually
called the `Klein--Gordon equation.'

Schr\"odinger's relativistic wave equation was derived by noting first
that, for a `Lorentz electron' of mass $m$ and charge $e$ in an external vector
potential $\mathbf{A}$ and Coulomb potential $\phi$, the Hamiltonian $H$ and momentum
$\mathbf{p}$ are related by\footnote{This is Lorentz invariant, because the quantities $\mathbf{A}$ and $\phi$ have the same Lorentz transformation property as $c\mathbf{p}$ and $E$. Schr\"odinger actually wrote $H$ and $\mathbf{p}$ in terms of partial derivatives of an action function, but this makes no difference to our present discussion.}
\begin{equation}
  0=(H+e\phi)^2-c^2(\mathbf{p}+e\mathbf{A}/c)^2-m^2c^4.
  \tag{1.1.2}\label{eq:1.1.2}
\end{equation}
For a free particle described by a plane wave
$\exp\{2\pi \ii(\boldsymbol{\kappa}\cdot\mathbf{x}-\nu t)\}$, the de
Broglie relations \eqref{eq:1.1.1} can be obtained by the identifications
\begin{equation}
  \mathbf{p}=h\boldsymbol{\kappa}\longrightarrow-\ii\hbar\boldsymbol{\nabla},
  \qquad
  E=h\nu\longrightarrow \ii\hbar\frac{\partial}{\partial t},
  \tag{1.1.3}\label{eq:1.1.3}
\end{equation}
where $\hbar$ is the convenient symbol (introduced later by Dirac) for $h/2\pi$.
By an admittedly formal analogy, Schr\"odinger guessed that an electron
in the external fields $\mathbf{A},\phi$ would be described by a wave function $\psi(\mathbf{x},t)$
satisfying the equation obtained by making the same replacements in
\eqref{eq:1.1.2}:
\begin{equation}
  0=\left[\left(\ii\hbar\frac{\partial}{\partial t}+e\phi\right)^2
  -c^2\left(-\ii\hbar\boldsymbol{\nabla}+\frac{e\mathbf{A}}{c}\right)^2
  -m^2c^4\right]\psi(\mathbf{x},t).
  \tag{1.1.4}\label{eq:1.1.4}
\end{equation}
In particular, for the stationary states of hydrogen we have $\mathbf{A}=0$ and
$\phi=e/4\pi r$, and $\psi$ has the time-dependence $\exp(-\ii Et/\hbar)$, so \eqref{eq:1.1.4} becomes
\begin{equation}
  0=\left[\left(E+\frac{e^2}{4\pi r}\right)^2
  -c^2\hbar^2\nabla^2-m^2c^4\right]\psi(\mathbf{x}).
  \tag{1.1.5}\label{eq:1.1.5}
\end{equation}
Solutions satisfying reasonable boundary conditions can be found for the
energy values~\hyperref[ref:1.9]{[9]}
\begin{equation}
  E=mc^2\left[1-\frac{\alpha^2}{2n^2}
  -\frac{\alpha^4}{2n^4}\left(\frac{n}{\ell+\tfrac12}-\frac34\right)
  +\cdots\right],
  \tag{1.1.6}\label{eq:1.1.6}
\end{equation}
where $\alpha=e^2/4\pi\hbar c$ is the `fine structure constant,' roughly $1/137$; $n$ is a
positive-definite integer, and $\ell$, the orbital angular momentum in units of
$\hbar$, is an integer with $0\leq\ell\leq n-1$. The $\alpha^2$ term gave good agreement
with the gross features of the hydrogen spectrum (the Lyman, Balmer,
etc. series) and, according to Dirac~\hyperref[ref:1.6]{[6]}, it was this agreement that led
Schr\"odinger eventually to develop his non-relativistic wave equation. On
the other hand, the $\alpha^4$ term gave a fine structure in disagreement with
existing accurate measurements of Friedrich Paschen~\hyperref[ref:1.10]{[10]}.

It is instructive here to compare Schr\"odinger's result with that of Arnold
Sommerfeld~\hyperref[ref:1.10c]{[10c]}, obtained using the rules of the old quantum theory:
\begin{equation}
  E=mc^2\left[1-\frac{\alpha^2}{2n^2}
  -\frac{\alpha^4}{2n^4}\left(\frac{n}{k}-\frac34\right)
  +\cdots\right].
  \tag{1.1.7}\label{eq:1.1.7}
\end{equation}
where $m$ is the electron mass. Here $k$ is an integer between 1 and $n$, which
in Sommerfeld's theory is given in terms of the orbital angular momentum
$\ell\hbar$ as $k=\ell+1$. This gave a fine structure splitting in agreement with
experiment: for instance, for $n=2$ Eq.~\eqref{eq:1.1.7} gives two levels ($k=1$
and $k=2$), split by the observed amount $\alpha^4mc^2/32$, or $4.53\cdot10^{-5}$ eV. In
contrast, Schr\"odinger's result \eqref{eq:1.1.6} gives an $n=2$ fine structure splitting
$\alpha^4mc^2/12$, considerably larger than observed.

Schr\"odinger correctly recognized that the source of this discrepancy
was his neglect of the spin of the electron. The splitting of atomic
energy levels by non-inverse-square electric fields in alkali atoms and by
weak external magnetic fields (the so-called anomalous Zeeman effect)
had revealed a multiplicity of states larger than could be accounted for
by the Bohr--Sommerfeld theory; this led George Uhlenbeck and Samuel
Goudsmit~\hyperref[ref:1.11]{[11]} in 1925 to suggest that the electron has an intrinsic angular
momentum $\hbar/2$. Also, the magnitude of the Zeeman splitting~\hyperref[ref:1.12]{[12]} allowed
them to estimate further that the electron has a magnetic moment
\begin{equation}
  \mu=\frac{e\hbar}{2mc}.
  \tag{1.1.8}\label{eq:1.1.8}
\end{equation}
It was clear that the electron's spin would be coupled to its orbital
angular momentum, so that Schr\"odinger's relativistic equation should not
be expected to give the correct fine structure splitting.

Indeed, by 1927 several authors~\hyperref[ref:1.13]{[13]} had been able to show that the
spin--orbit coupling was able to account for the discrepancy between Schr\"odinger's result \eqref{eq:1.1.6} and experiment. There are really two effects here:
one is a direct coupling between the magnetic moment \eqref{eq:1.1.8} and the
magnetic field felt by the electron as it moves through the electrostatic
field of the atom; the other is the relativistic `Thomas precession' caused
(even in the absence of a magnetic moment) by the circular motion of
the spinning electron~\hyperref[ref:1.14]{[14]}. Together, these two effects were found to lift the
level with total angular momentum $j=\ell+\tfrac12$ to the energy \eqref{eq:1.1.7} given
by Sommerfeld for $k=\ell+1=j+\tfrac12$, while the level with $j=\ell-\tfrac12$ was
lowered to the value given by Sommerfeld for $k=\ell=j+\tfrac12$. Thus the
energy was found to depend only on $n$ and $j$, but not separately on $\ell$:
\begin{equation}
  E=mc^2\left[1-\frac{\alpha^2}{2n^2}
  -\frac{\alpha^4}{2n^4}\left(\frac{n}{j+\tfrac12}-\frac34\right)
  +\cdots\right].
  \tag{1.1.9}\label{eq:1.1.9}
\end{equation}
By accident Sommerfeld's theory had given the correct magnitude of the
splitting in hydrogen ($j+\tfrac12$ like $k$ runs over integer values from 1 to $n$)
though it was wrong as to the assignment of orbital angular momentum
values $\ell$ to these various levels. In addition, the multiplicity of the fine
structure levels in hydrogen was now predicted to be 2 for $j=\tfrac12$ and
$2(2j+1)$ for $j>\tfrac12$ (corresponding to $\ell$ values $j\pm\tfrac12$), in agreement with
experiment.

Despite these successes, there still was not a thorough relativistic theory
which incorporated the electron's spin from the beginning. Such a theory
was discovered in 1928 by Paul Dirac. However, he did not set out
simply to make a relativistic theory of the spinning electron; instead, he
approached the problem by posing a question that would today seem
very strange. At the beginning of his 1928 paper~\hyperref[ref:1.15]{[15]}, he asks `why Nature
should have chosen this particular model for the electron, instead of
being satisfied with the point charge.' To us today, this question is like
asking why bacteria have only one cell: having spin $\hbar/2$ is just one of
the properties that define a particle as an electron, rather than one of the
many other types of particles with various spins that are known today.
However, in 1928 it was possible to believe that all matter consisted
of electrons, and perhaps something similar with positive charge in the
atomic nucleus. Thus, in the spirit of the times in which it was asked,
Dirac's question can be restated: `Why do the fundamental constituents
of matter have to have spin $\hbar/2$?'

For Dirac, the key to this question was the requirement that probabilities
must be positive. It was known~\hyperref[ref:1.16]{[16]} that the probability density for the non-relativistic Schr\"odinger equation is $|\psi|^2$, and that this satisfies a continuity
equation of the form
\begin{equation*}
  \frac{\partial}{\partial t}|\psi|^2
  -\frac{\ii\hbar}{2m}\boldsymbol{\nabla}\cdot
  (\psi^*\boldsymbol{\nabla}\psi-\psi\boldsymbol{\nabla}\psi^*)=0
\end{equation*}
so the space-integral of $|\psi|^2$ is time-independent. On the other hand,
the only probability density $\rho$ and current $\mathbf{J}$, which can be formed from
solutions of the relativistic Schr\"odinger equation and which satisfy a
conservation law,
\begin{equation}
  \frac{\partial\rho}{\partial t}+\boldsymbol{\nabla}\cdot\mathbf{J}=0,
  \tag{1.1.10}\label{eq:1.1.10}
\end{equation}
are of the form
\begin{align}
  \rho&=N\operatorname{Im}\psi^*
  \left(\frac{\partial}{\partial t}-\frac{ie\phi}{\hbar}\right)\psi,
  \tag{1.1.11}\label{eq:1.1.11}\\
  \mathbf{J}&=Nc^2\operatorname{Im}\psi^*
  \left(\boldsymbol{\nabla}+\frac{ie\mathbf{A}}{\hbar c}\right)\psi,
  \tag{1.1.12}\label{eq:1.1.12}
\end{align}
with $N$ an arbitrary constant. It is not possible to identify $\rho$ as the
probability density, because (with or without an external potential $\phi$) $\rho$
does not have definite sign. To quote Dirac's reminiscences~\hyperref[ref:1.17]{[17]} about this
problem
\begin{quote}
I remember once when I was in Copenhagen, that Bohr
asked me what I was working on and I told him I was trying
to get a satisfactory relativistic theory of the electron, and Bohr
said `But Klein and Gordon have already done that!' That
answer first rather disturbed me. Bohr seemed quite satisfied
by Klein's solution, but I was not because of the negative
probabilities that it led to. I just kept on with it, worrying about
getting a theory which would have only positive probabilities.
\end{quote}

According to George Gamow~\hyperref[ref:1.18]{[18]}, Dirac found the answer to this problem
on an evening in 1928 while staring into a fireplace at St John's College,
Cambridge. He realized that the reason that the Klein--Gordon (or
relativistic Schr\"odinger) equation can give negative probabilities is that
the $\rho$ in the conservation equation \eqref{eq:1.1.10} involves a time-derivative of the
wave function. This in turn happens because the wave function satisfies
a differential equation of second order in the time. The problem therefore
was to replace this wave equation with another one of first order in time
derivatives, like the non-relativistic Schr\"odinger equation.

Suppose the electron wave function is a multi-component quantity $\psi_n(\mathbf{x})$,
which satisfies a wave equation of the form,
\begin{equation}
  \ii\hbar\frac{\partial\psi}{\partial t}=\mathcal{H}\psi,
  \tag{1.1.13}\label{eq:1.1.13}
\end{equation}
where $\mathcal{H}$ is some matrix function of space derivatives. In order to have a
chance at a Lorentz-invariant theory, we must suppose that because the
equation is linear in time-derivatives, it is also linear in space-derivatives,
so that $\mathcal{H}$ takes the form:
\begin{equation}
  \mathcal{H}=-\ii\hbar c\boldsymbol{\alpha}\cdot\boldsymbol{\nabla}+\beta mc^2,
  \tag{1.1.14}\label{eq:1.1.14}
\end{equation}
where $\alpha_1$, $\alpha_2$, $\alpha_3$, and $\beta$ are constant matrices. From \eqref{eq:1.1.13} we can derive
the second-order equation
\begin{align*}
  -\hbar^2\frac{\partial^2\psi}{\partial t^2}
  =\mathcal{H}^2\psi
  ={}&-\hbar^2c^2\alpha_i\alpha_j
  \frac{\partial^2\psi}{\partial x_i\partial x_j}\\
  &-\ii\hbar mc^3(\alpha_i\beta+\beta\alpha_i)
  \frac{\partial\psi}{\partial x_i}
  +m^2c^4\beta^2\psi.
\end{align*}
(The summation convention is in force here; $i$ and $j$ run over the values
1, 2, 3, or $x,y,z$.) But this must agree with the free-field form of the
relativistic Schr\"odinger equation \eqref{eq:1.1.4}, which just expresses the relativistic
relation between momentum and energy. Therefore, the matrices $\boldsymbol{\alpha}$ and $\beta$
must satisfy the relations
\begin{align}
  \alpha_i\alpha_j+\alpha_j\alpha_i&=2\delta_{ij}\mathbf{1},
  \tag{1.1.15}\label{eq:1.1.15}\\
  \alpha_i\beta+\beta\alpha_i&=0,
  \tag{1.1.16}\label{eq:1.1.16}\\
  \beta^2&=\mathbf{1},
  \tag{1.1.17}\label{eq:1.1.17}
\end{align}
where $\delta_{ij}$ is the Kronecker delta (unity for $i=j$; zero for $i\neq j$) and $\mathbf{1}$ is
the unit matrix. Dirac found a set of $4\times4$ matrices which satisfy these
relations
\begin{equation}
\begin{gathered}
  \alpha_1=
  \begin{bmatrix}
    0&0&0&1\\
    0&0&1&0\\
    0&1&0&0\\
    1&0&0&0
  \end{bmatrix},
  \qquad
  \alpha_2=
  \begin{bmatrix}
    0&0&0&-\ii\\
    0&0&\ii&0\\
    0&-\ii&0&0\\
    i&0&0&0
  \end{bmatrix},\\[1ex]
  \alpha_3=
  \begin{bmatrix}
    0&0&1&0\\
    0&0&0&-1\\
    1&0&0&0\\
    0&-1&0&0
  \end{bmatrix},
  \qquad
  \beta=
  \begin{bmatrix}
    1&0&0&0\\
    0&1&0&0\\
    0&0&-1&0\\
    0&0&0&-1
  \end{bmatrix}.
\end{gathered}
\tag{1.1.18}\label{eq:1.1.18}
\end{equation}

To show that this formalism is Lorentz-invariant, Dirac multiplied
Eq.~\eqref{eq:1.1.13} on the left with $\beta$, so that it could be put in the form
\begin{equation}
  \left[\hbar c\gamma^\mu\frac{\partial}{\partial x^\mu}+mc^2\right]\psi=0,
  \tag{1.1.19}\label{eq:1.1.19}
\end{equation}
where
\begin{equation}
  \gamma^i\equiv-\ii\beta\alpha_i,
  \qquad
  \gamma^0\equiv-\ii\beta.
  \tag{1.1.20}\label{eq:1.1.20}
\end{equation}
(The Greek indices $\mu$, $\nu$, etc. will now run over the values $0,1,2,3$, with
$x^0=ct$. Dirac used $x_4=ict$, and correspondingly $\gamma_4=\beta$.) The matrices
$\gamma^\mu$ satisfy the anticommutation relations
\begin{equation}
  \frac12(\gamma^\mu\gamma^\nu+\gamma^\nu\gamma^\mu)
  =\eta^{\mu\nu}\equiv
  \begin{cases}
    +1,&\mu=\nu=1,2,3,\\
    -1,&\mu=\nu=0,\\
    0,&\mu\neq\nu.
  \end{cases}
  \tag{1.1.21}\label{eq:1.1.21}
\end{equation}
Dirac noted that these anticommutation relations are Lorentz-invariant,
in the sense that they are also satisfied by the matrices
$\Lambda^\mu{}_{\nu}\gamma^\nu$, where $\Lambda$ is
any Lorentz transformation. He concluded from this that $\Lambda^\mu{}_{\nu}\gamma^\nu$ must be
related to $\gamma^\mu$ by a similarity transformation:
\begin{equation*}
  \Lambda^\mu{}_{\nu}\gamma^\nu=S^{-1}(\Lambda)\gamma^\mu S(\Lambda).
\end{equation*}
It follows that the wave equation is invariant if, under a Lorentz transformation
$x^\mu\to\Lambda^\mu{}_{\nu}x^\nu$, the wave function undergoes the matrix transformation
$\psi\to S(\Lambda)\psi$. (These matters are discussed more fully, from a rather
different point of view, in Chapter 5.)

To study the behavior of electrons in an arbitrary external electromagnetic
field, Dirac followed the `usual procedure' of making the replacements
\begin{equation}
  \ii\hbar\frac{\partial}{\partial t}
  \longrightarrow \ii\hbar\frac{\partial}{\partial t}+e\phi,
  \qquad
  -\ii\hbar\boldsymbol{\nabla}\longrightarrow-\ii\hbar\boldsymbol{\nabla}+\frac{e}{c}\mathbf{A}
  \tag{1.1.22}\label{eq:1.1.22}
\end{equation}
as in Eq.~\eqref{eq:1.1.4}. The wave equation \eqref{eq:1.1.13} then takes the form
\begin{equation}
  \left(\ii\hbar\frac{\partial}{\partial t}+e\phi\right)\psi
  =(-\ii\hbar c\boldsymbol{\nabla}+e\mathbf{A})\cdot\boldsymbol{\alpha}\psi+mc^2\beta\psi.
  \tag{1.1.23}\label{eq:1.1.23}
\end{equation}
Dirac used this equation to show that in a central field, the conservation
of angular momentum takes the form
\begin{equation}
  [\mathcal{H},-\ii\hbar\mathbf{r}\cdot\boldsymbol{\nabla}+\hbar\boldsymbol{\sigma}/2]=0,
  \tag{1.1.24}\label{eq:1.1.24}
\end{equation}
where $\mathcal{H}$ is the matrix differential operator \eqref{eq:1.1.14} and $\boldsymbol{\sigma}$ is the $4\times4$
version of the spin matrix introduced earlier by Pauli~\hyperref[ref:1.19]{[19]}
\begin{equation}
  \boldsymbol{\sigma}=
  \begin{bmatrix}
    0&0&1&0\\
    0&0&0&1\\
    1&0&0&0\\
    0&1&0&0
  \end{bmatrix}\boldsymbol{\alpha}.
  \tag{1.1.25}\label{eq:1.1.25}
\end{equation}
Since each component of $\boldsymbol{\sigma}$ has eigenvalues $\pm1$, the presence of the extra
term in \eqref{eq:1.1.24} shows that the electron has intrinsic angular momentum
$\hbar/2$.

Dirac also iterated Eq.~\eqref{eq:1.1.23}, obtaining a second-order equation,
which turned out to have just the same form as the Klein--Gordon equation
\eqref{eq:1.1.4} except for the presence on the right-hand-side of two additional
terms
\begin{equation}
  [-e\hbar c\boldsymbol{\sigma}\cdot\mathbf{B}-ie\hbar c\boldsymbol{\alpha}\cdot\mathbf{E}]\psi.
  \tag{1.1.26}\label{eq:1.1.26}
\end{equation}
For a slowly moving electron, the first term dominates, and represents a
magnetic moment in agreement with the value \eqref{eq:1.1.8} found by Goudsmit
and Uhlenbeck~\hyperref[ref:1.11]{[11]}. As Dirac recognized, this magnetic moment, together
with the relativistic nature of the theory, guaranteed that this theory
would give a fine structure splitting in agreement (to order $\alpha^4mc^2$) with
that found by Heisenberg, Jordan, and Charles G. Darwin~\hyperref[ref:1.13]{[13]}. A little later,
an `exact' formula for the hydrogen energy levels in Dirac's theory was
derived by Darwin~\hyperref[ref:1.20]{[20]} and Gordon~\hyperref[ref:1.21]{[21]}
\begin{equation}
  E=mc^2\left(
  1+\frac{\alpha^2}{
  \left\{n-j-\tfrac12+
  \left[(j+\tfrac12)^2-\alpha^2\right]^{1/2}\right\}^2}
  \right)^{-1/2}.
  \tag{1.1.27}\label{eq:1.1.27}
\end{equation}
The first three terms of a power series expansion in $\alpha^2$ agree with the
approximate result \eqref{eq:1.1.9}.

This theory achieved Dirac's primary aim: a relativistic formalism with
positive probabilities. From \eqref{eq:1.1.13} we can derive a continuity equation
\begin{equation}
  \frac{\partial\rho}{\partial t}+\boldsymbol{\nabla}\cdot\mathbf{J}=0
  \tag{1.1.28}\label{eq:1.1.28}
\end{equation}
with
\begin{equation}
  \rho=|\psi|^2,
  \qquad
  \mathbf{J}=c\psi^\dagger\boldsymbol{\alpha}\psi,
  \tag{1.1.29}\label{eq:1.1.29}
\end{equation}
so that the positive quantity $|\psi|^2$ can be interpreted as a probability
density, with constant total probability $\int|\psi|^2\dd^3x$. However, there was
another difficulty which Dirac was not immediately able to resolve.

For a given momentum $\mathbf{p}$, the wave equation \eqref{eq:1.1.13} has four solutions
of the plane wave form
\begin{equation}
  \psi\propto\exp\left[\frac{\ii}{\hbar}(\mathbf{p}\cdot\mathbf{x}-Et)\right].
  \tag{1.1.30}\label{eq:1.1.30}
\end{equation}
Two solutions with $E=+\sqrt{\mathbf{p}^2c^2+m^2c^4}$ correspond to the two spin states
of an electron with $J_z=\pm\hbar/2$. The other two solutions have
$E=-\sqrt{\mathbf{p}^2c^2+m^2c^4}$, and no obvious physical interpretation. As Dirac pointed
out, this problem arises also for the relativistic Schr\"odinger equation: for
each $\mathbf{p}$, there are two solutions of the form \eqref{eq:1.1.30}, one with positive $E$
and one with negative $E$.

Of course, even in classical physics, the relativistic relation
$E^2=\mathbf{p}^2c^2+m^2c^4$ has two solutions, $E=\pm\sqrt{\mathbf{p}^2c^2+m^2c^4}$. However,
in classical physics we can simply assume that the only physical particles
are those with positive $E$. Since the positive solutions have $E>mc^2$ and
the negative ones have $E<-mc^2$, there is a finite gap between them,
and no continuous process can take a particle from positive to negative
energy.

The problem of negative energies is much more troublesome in relativistic
quantum mechanics. As Dirac pointed out in his 1928 paper~\hyperref[ref:1.15]{[15]}, the
interaction of electrons with radiation can produce transitions in which a
positive-energy electron falls into a negative-energy state, with the energy
carried off by two or more photons. Why then is matter stable?

In 1930 Dirac offered a remarkable solution~\hyperref[ref:1.22]{[22]}. Dirac's proposal was
based on the exclusion principle, so a few words about the history of this
principle are in order here.

The periodic table of the elements and the systematics of X-ray spectroscopy
had together by 1924 revealed a pattern in the population of
atomic energy levels by electrons~\hyperref[ref:1.23]{[23]}. The maximum number $N_n$ of electrons
in a shell characterized by principal quantum number $n$ is given by twice
the number of orbital states with that $n$
\begin{equation}
  N_n=2\sum_{\ell=0}^{n-1}(2\ell+1)=2n^2=2,8,18,\ldots.
  \tag{1.1.31}\label{eq:1.1.31}
\end{equation}
Wolfgang Pauli~\hyperref[ref:1.24]{[24]} in 1925 suggested that this pattern could be understood if
$N_n$ is the total number of possible states in the $n$th shell, and if in addition
there is some mysterious `exclusion principle' which forbids more than one
electron from occupying the same state. He explained the puzzling factor
2 in \eqref{eq:1.1.31} as due to a `peculiar, classically non-describable duplexity' of
the electron states, and as we have seen this was understood a little later
as due to the spin of the electron~\hyperref[ref:1.11]{[11]}. The exclusion principle answered a
question that had remained obscure in the old atomic theory of Bohr and
Sommerfeld: why do not all the electrons in heavy atoms fall down into
the shell of lowest energy? Subsequently Pauli's exclusion principle was
formalized by a number of authors~\hyperref[ref:1.25]{[25]} as the requirement that the wave
function of a multi-electron system is antisymmetric in the coordinates,
orbital and spin, of all the electrons. This principle was incorporated into
statistical mechanics by Enrico Fermi~\hyperref[ref:1.26]{[26]} and Dirac~\hyperref[ref:1.27]{[27]}, and for this reason
particles obeying the exclusion principle are generally called `fermions,'
just as particles like photons for which the wave function is symmetric
and which obey the statistics of Bose and Einstein are called `bosons.' The
exclusion principle has played a fundamental role in the theory of metals,
white dwarf and neutron stars, etc., as well as in chemistry and atomic
physics, but a discussion of these matters would take us too far afield
here.

Dirac's proposal was that the positive energy electrons cannot fall down
into negative energy states because `all the states of negative energy are
occupied except perhaps a few of small velocity.' The few vacant states,
or `holes,' in the sea of negative energy electrons behave like particles with
opposite quantum numbers: positive energy and positive charge. The only
particle with positive charge that was known at that time was the proton,
and as Dirac later recalled~\hyperref[ref:1.27a]{[27a]}, `the whole climate of opinion at that time
was against new particles' so Dirac identified his holes as protons; in fact,
the title of his 1930 article~\hyperref[ref:1.22]{[22]} was `A Theory of Electrons and Protons.'

The hole theory faced a number of immediate difficulties. One obvious
problem was raised by the infinite charge density of the ubiquitous
negative-energy electrons: where is their electric field? Dirac proposed to
reinterpret the charge density appearing in Maxwell's equations as `the
departure from the normal state of electrification of the world.' Another
problem has to do with the huge dissimilarity between the observed
masses and interactions of the electrons and protons. Dirac hoped that
Coulomb interactions between electrons would somehow account for these
differences but Hermann Weyl~\hyperref[ref:1.28]{[28]} showed that the hole theory was in fact
entirely symmetric between negative and positive charge. Finally, Dirac~\hyperref[ref:1.22]{[22]}
predicted the existence of an electron--proton annihilation process in which
a positive-energy electron meets a hole in the sea of negative-energy electrons
and falls down into the unoccupied level, emitting a pair of gamma
ray photons. By itself this would not have created difficulties for the hole
theory; it was even hoped by some that this would provide an explanation,
then lacking, of the energy source of the stars. However, it was
soon pointed out~\hyperref[ref:1.29]{[29]} by Julius Robert Oppenheimer and Igor Tamm that
electron--proton annihilation in atoms would take place at much too fast
a rate to be consistent with the observed stability of ordinary matter. For
these reasons, by 1931 Dirac had changed his mind, and decided that the
holes would have to appear not as protons but as a new sort of positively
charged particle, of the same mass as the electron~\hyperref[ref:1.29a]{[29a]}.

The second and third of these problems were eliminated by the discovery
of the positron by Carl D. Anderson~\hyperref[ref:1.30]{[30]}, who apparently did not know of
this prediction by Dirac. On August 2, 1932, a peculiar cosmic ray track
was observed in a Wilson cloud chamber subjected to a 15 kG magnetic
field. The track was observed to curve in a direction that would be
expected for a positively charged particle, and yet its range was at least
ten times greater than the expected range of a proton! Both the range
and the specific ionization of the track were consistent with the hypothesis
that this was a new particle which differs from the electron only in the
sign of its charge, as would be expected for one of Dirac's holes. (This
discovery had been made earlier by P.M.S. Blackett, but not immediately
published by him. Anderson quotes press reports of evidence for light
positive particles in cosmic ray tracks, obtained by Blackett and Giuseppe
Occhialini.) Thus it appeared that Dirac was wrong only in his original
identification of the hole with the proton.

The discovery of the more-or-less predicted positron, together with the
earlier successes of the Dirac equation in accounting for the magnetic
moment of the electron and the fine structure of hydrogen, gave Dirac's
theory a prestige that it has held for over six decades. However, although
there seems little doubt that Dirac's theory will survive in some form in
any future physical theory, there are serious reasons for being dissatisfied
with its original rationale:

\medskip
\noindent
(i)\quad Dirac's analysis of the problem of negative probabilities in Schr\"odinger's
relativistic wave equation would seem to rule out the existence
of any particle of zero spin. Yet even in the 1920s particles of zero spin
were known---for instance, the hydrogen atom in its ground state, and
the helium nucleus. Of course, it could be argued that hydrogen atoms
and alpha particles are not elementary, and therefore do not need to
be described by a relativistic wave equation, but it was not (and still is
not) clear how the idea of elementarity is incorporated in the formalism
of relativistic quantum mechanics. Today we know of a large number
of spin zero particles---$\pi$ mesons, $K$ mesons, and so on---that are
no less elementary than the proton and neutron. We also know of spin
one particles---the $W^\pm$ and $Z^0$---which seem as elementary as the
electron or any other particle. Further, apart from effects of the strong
interactions, we would today calculate the fine structure of `mesonic
atoms,' consisting of a spinless negative $\pi$ or $K$ meson bound to an
atomic nucleus, from the stationary solutions of the relativistic Klein--Gordon--Schr\"odinger equation! Thus, it is difficult to agree that there is
anything fundamentally wrong with the relativistic equation for zero spin
that \emph{forced} the development of the Dirac equation---the problem simply
is that the electron happens to have spin $\hbar/2$, not zero.

\medskip
\noindent
(ii)\quad As far as we now know, for \emph{every} kind of particle there is an
`antiparticle' with the same mass and opposite charge. (Some purely
neutral particles, such as the photon, are their own antiparticles.) But
how can we interpret the antiparticles of charged \emph{bosons}, such as the
$\pi^\pm$ mesons or $W^\pm$ particles, as holes in a sea of negative energy states?
For particles quantized according to the rules of Bose--Einstein statistics,
there is no exclusion principle, and hence nothing to keep positive-energy
particles from falling down into the negative-energy states, occupied or
not. And if the hole theory does not work for bosonic antiparticles, why
should we believe it for fermions? I asked Dirac in 1972 how he then felt
about this point; he told me that he did not regard bosons like the pion
or $W^\pm$ as `important.' In a lecture~\hyperref[ref:1.27a]{[27a]} a few years later, Dirac referred
to the fact that for bosons `we no longer have the picture of a vacuum
with negative energy states filled up', and remarked that in this case `the
whole theory becomes more complicated.' The next section will show
how the development of quantum field theory made the interpretation of
antiparticles as holes unnecessary, even though unfortunately it lingers
on in many textbooks. To quote Julian Schwinger~\hyperref[ref:1.30a]{[30a]}, `The picture of an
infinite sea of negative energy electrons is now best regarded as a historical
curiosity, and forgotten.'

\medskip
\noindent
(iii)\quad One of the great successes of the Dirac theory was its correct
prediction of the magnetic moment of the electron. This was particularly
striking, as the magnetic moment \eqref{eq:1.1.8} is twice as large as would be
expected for the orbital motion of a charged point particle with angular
momentum $\hbar/2$; this factor of 2 had remained mysterious until Dirac's
theory. However, there is really nothing in Dirac's line of argument that
leads unequivocally to this particular value for the magnetic moment. At
the point where we brought electric and magnetic fields into the wave
equation \eqref{eq:1.1.23}, we could just as well have added a `Pauli term'~\hyperref[ref:1.31]{[31]}
\begin{equation}
  \kappa c\beta[\gamma^\mu,\gamma^\nu]\psi F_{\mu\nu}
  \tag{1.1.32}\label{eq:1.1.32}
\end{equation}
with arbitrary coefficient $\kappa$. (Here $F_{\mu\nu}$ is the usual electromagnetic field
strength tensor, with $F^{12}=B_3$, $F^{01}=E_1$, etc.) This term could be
obtained by first adding a term to the free-field equations proportional
to $[\gamma^\mu,\gamma^\nu](\partial^2/\partial x^\mu\partial x^\nu)\psi$, which of course equals zero, and then making
the substitutions \eqref{eq:1.1.22} as before. A more modern approach would be
simply to remark that the term \eqref{eq:1.1.32} is consistent with all accepted
invariance principles, including Lorentz invariance and gauge invariance,
and so there is no reason why such a term should not be included in the
field equations. (See Section 12.3.) This term would give an additional
contribution proportional to $\kappa$ to the magnetic moment of the electron, so
apart from the possible demand for a purely formal simplicity, there was
no reason to expect any particular value for the magnetic moment of the
electron in Dirac's theory.

As we shall see in this book, these problems were all eventually to be
solved (or at least clarified) through the development of quantum field
theory.
